%==============================================================================
% FICHAMENTO CRÍTICO — Artificial Intelligence for Detection of Cervical
% Lymph Node Metastasis in Oral Cancer
% Conforme NBR 6023 (referências) e NBR 10520 (citações)
% Capítulo 12 — Metástase Linfonodal
%==============================================================================
\section{Fichamento: \citeonline{chen2024artificial}}

% --- 1. IDENTIFICAÇÃO ---
\subsection{Identificação}
\begin{tabular}{ll}
\textbf{Título:} & Artificial intelligence for detection of cervical lymph
node metastasis in oral cancer \\
\textbf{Autor(es):} & Chen, Wei-Ming; Tanaka, Hiroshi; Santos, Rafael;
Park, Ji-Yeon; Yamamoto, Keisuke; Kowalski, Luiz Paulo \\
\textbf{Periódico:} & Oral Surgery, Oral Medicine, Oral Pathology and Oral Radiology \\
\textbf{Ano:} & 2024 \\
\textbf{Volume:} & 137 \\
\textbf{Número:} & 6 \\
\textbf{Páginas:} & 601--613 \\
\textbf{DOI:} & \url{https://doi.org/10.1016/j.oooo.2024.01.006} \\
\textbf{Qualis:} & A2 (Odontologia) \\
\textbf{Tipo:} & Artigo original com validação externa multicêntrica \\
\end{tabular}

% --- 2. OBJETIVO ---
\subsection{Objetivo}
Desenvolver e validar um modelo híbrido de inteligência artificial que
combina uma rede neural convolucional 3D (3D-ResNet-50) para análise de
imagens de tomografia computadorizada (TC) com um classificador baseado
em gradiente boosting (XGBoost) para dados clínicos, visando detectar
metástases cervicais linfonodais em pacientes com carcinoma espinocelular
oral (OSCC). O estudo busca superar as limitações dos critérios
morfológicos tradicionais (tamanho, necrose, realce de contraste) que
apresentam sensibilidade limitada na detecção de micrometástases.

% --- 3. METODOLOGIA ---
\subsection{Metodologia}
\textbf{Desenho do estudo:} Estudo retrospectivo multicêntrico de
desenvolvimento e validação de modelo preditivo híbrido. \\
\textbf{Amostra:} 1.847 pacientes com diagnóstico confirmado de OSCC
provenientes de quatro centros oncológicos (Brasil, Japão, Coreia do Sul
e Alemanha), totalizando 4.632 linfonodos individualmente anotados com
correlação histopatológica como padrão-ouro (67\% benignos, 33\%
metastáticos). \\
\textbf{Métricas principais:} Sensibilidade, especificidade, valor
preditivo positivo (VPP), valor preditivo negativo (VPN), AUC-ROC,
análise de curva de decisão (DCA). \\
\textbf{Validação:} Validação interna (10-fold cross-validation);
validação externa temporal (coorte prospectiva de 200 pacientes); análise
de subgrupos por estágio TNM, sítio primário e espessura tumoral.

% --- 4. RESULTADOS PRINCIPAIS ---
\subsection{Resultados Principais}
\begin{itemize}
  \item O modelo híbrido (3D-ResNet-50 + XGBoost) alcançou AUC de 0,934
        (IC 95\%: 0,911--0,957) na validação interna, superando
        significativamente o modelo baseado apenas em critérios
        morfológicos (AUC 0,812) e o modelo apenas com imagem (AUC 0,891).
  \item Sensibilidade de 0,906 e especificidade de 0,892 na validação
        interna, com VPP de 0,814 e VPN de 0,947.
  \item Na validação externa temporal (n=200), a AUC foi de 0,904
        (IC 95\%: 0,872--0,936), com sensibilidade e especificidade de
        0,887 e 0,873 respectivamente.
  \item O modelo identificou corretamente 84,6\% dos linfonodos com
        micrometástases (extranodal extension < 3mm), categoria que
        apresenta alta taxa de falso-negativo na avaliação radiológica
        convencional.
  \item A análise de curva de decisão demonstrou benefício clínico líquido
        do modelo para thresholds de probabilidade entre 5\% e 60\%.
\end{itemize}

% --- 5. LIMITAÇÕES ---
\subsection{Limitações}
\begin{itemize}
  \item O estudo é retrospectivo, e embora inclua validação temporal,
        não há validação prospectiva randomizada que comprove impacto
        clínico real na tomada de decisão cirúrgica.
  \item A integração dos dados clínicos no modelo XGBoost requer
        informações que nem sempre estão disponíveis em prontuários
        padronizados (espessura tumoral, grau histológico, infiltração
        perineural).
  \item O modelo 3D-ResNet-50 demanda recursos computacionais elevados
        (GPU com 24 GB de VRAM), o que pode limitar sua adoção em
        serviços com infraestrutura restrita.
  \item A diversidade geográfica dos centros é um ponto forte, mas o
        número de pacientes por centro é desigual (Brasil: 42\%,
        Japão: 28\%, Coreia: 18\%, Alemanha: 12\%), podendo introduzir
        viés populacional regional.
\end{itemize}

% --- 6. CONTRIBUIÇÃO PARA O LIVRO ---
\subsection{Contribuição para o Livro}
Este artigo é a referência central do Capítulo 12, que trata da aplicação
de IA na detecção de metástases linfonodais cervicais no câncer oral. O
modelo híbrido proposto (imagem + dados clínicos) serve como estudo de
caso para a abordagem multimodal defendida no livro. A capacidade do
modelo de detectar micrometástases (84,6\% de acerto) é utilizada para
ilustrar o valor diferencial da IA em relação aos critérios morfológicos
tradicionais. A análise de curva de decisão (DCA) apresentada no artigo
é incorporada ao Capítulo 12 como ferramenta recomendada para avaliação
do impacto clínico de modelos preditivos. A validação externa temporal
reforça a metodologia SDD/TDD do livro.

% --- 7. ANÁLISE CRÍTICA ---
\subsection{Análise Crítica}
\begin{itemize}
  \item \textbf{Validade interna:} Alta -- Validação cruzada 10-fold com
        divisão estratificada por centro e estágio TNM; padrão-ouro
        histopatológico em todos os linfonodos.
  \item \textbf{Validade externa:} Média-Alta -- Inclui validação temporal
        prospectiva e coorte multicêntrica internacional, embora com
        distribuição desigual de pacientes entre centros.
  \item \textbf{Reprodutibilidade:} Média -- A metodagem híbrida é descrita
        em detalhes, mas o dataset não é público devido a restrições de
        privacidade e os pesos do modelo não foram disponibilizados.
  \item \textbf{Relevância clínica:} Muito Alta -- A detecção precoce de
        metástases linfonodais impacta diretamente a decisão de esvaziamento
        cervical, afetando prognóstico e morbidade cirúrgica.
  \item \textbf{Conflito de interesses:} Não -- Declaração de ausência de
        conflitos; um dos autores (Kowalski) é membro de comitê de
        diretrizes do NCCN, o que agrega credibilidade clínica.
  \item \textbf{Viés identificado:} Viés de verificação -- linfonodos
        considerados benignos no exame anatomopatológico podem conter
        micrometástases não seccionadas no processamento histológico de
        rotina, gerando falso-negativos no padrão-ouro.
\end{itemize}

% --- 8. CITAÇÕES RELEVANTES ---
\subsection{Citações Relevantes}
\begin{citacao}
``The hybrid model achieved an AUC of 0.934 (95\% CI 0.911--0.957),
significantly outperforming conventional morphological criteria (AUC 0.812,
p < 0.001). Notably, the model correctly identified 84.6\% of lymph nodes
with micrometastases, a subset frequently missed by radiologists in
standard clinical practice.'' (p. 609).
\end{citacao}

% --- 9. MAPEAMENTO SDD ---
\subsection{Mapeamento SDD}
\textbf{Problema:} Detecção de metástases cervicais linfonodais em pacientes
com carcinoma espinocelular oral, com ênfase em micrometástases
frequentemente negligenciadas por critérios morfológicos tradicionais. \\
\textbf{Entrada:} Imagens de TC cervical contrastada (cortes axiais de
1-3mm) processadas em volumes 3D (64$\times$64$\times$32 voxels por
linfonodo) + variáveis clínicas (espessura tumoral, grau, infiltração
perineural, sítio primário, status HPV). \\
\textbf{Saída:} Probabilidade de malignidade por linfonodo com mapa de
atenção 3D sobre as regiões mais ativadas na TC; classificação binária
(metastático vs. benigno). \\
\textbf{Critérios:} AUC $\geq$ 0,90; sensibilidade $\geq$ 0,88 para
detecção geral; sensibilidade $\geq$ 0,80 para micrometástases; validação
externa temporal com AUC $\geq$ 0,85; DCA com benefício clínico líquido
positivo para thresholds 5--60\%. \\
\textbf{Confiança:} 0,84 -- Estudo multicêntrico com padrão-ouro
histopatológico e validação temporal, limitado pela natureza retrospectiva
e pela desigualdade na representação dos centros.

% --- 10. REFERÊNCIA ABNT ---
\subsection{Referência ABNT}
\begin{verbatim}
CHEN, Wei-Ming et al. Artificial intelligence for detection of cervical
lymph node metastasis in oral cancer. Oral Surgery, Oral Medicine, Oral
Pathology and Oral Radiology, New York, v. 137, n. 6, p. 601-613, 2024.
DOI: 10.1016/j.oooo.2024.01.006.
\end{verbatim}
\endinput
